\documentclass[11pt]{article}
% commande d'exécution simple :
% pandoc RAPPORT.md --template=template_jpp12.tex -o RAPPORT.pdf --pdf-engine=xelatex
% pandoc RAPPORT.md --template=template_jpp12.tex --listings -o RAPPORT.pdf --pdf-engine=xelatex
% pandoc Rapport.md --template=template_jpp12.tex --listings --citeproc -o rapport.pdf --pdf-engine=xelatex
% pour avoir les emojis facilement, utiliser lualatex
% pandoc RAPPORT.md --template=template_jpp12.tex --listings -o RAPPORT.pdf --pdf-engine=lualatex
% ==================== ENCODAGE ET LANGUE ====================
\usepackage[T1]{fontenc}
\usepackage[utf8]{inputenc}
\usepackage[french]{babel}
\usepackage{fontspec}
\usepackage{xeCJK}
\frenchbsetup{StandardLists=true}
% ==================== POLICES PRINCIPALES ====================
\setmainfont{Libertinus Serif}
\setmonofont{Libertinus Mono}
\setCJKmainfont{Noto Serif CJK SC}
% ==================== EMOJIS COULEUR ====================
%\usepackage[unicode]{emoji}   % Active la saisie directe des emojis (✅, 👍, …)
%\setemojifont{Noto Color Emoji} % Police recommandée (présente sur Linux, Windows, macOS)
% ==================== MATHÉMATIQUES ====================
\usepackage{amsmath}
\usepackage{mathtools}
\usepackage{amssymb}
\usepackage{unicode-math}
\setmathfont{Libertinus Math}
\usepackage{dsfont}
% ==================== COMMANDES PERSONNALISÉES ====================
\usepackage{pifont}
\newcommand{\bra}[1]{\langle #1|}
\newcommand{\ket}[1]{|#1\rangle}
\newcommand{\braket}[2]{\langle #1|#2\rangle}
\newcommand{\R}{\mathbb{R}}
\newcommand{\C}{\mathbb{C}}
\newcommand{\N}{\mathbb{N}}
\newcommand{\Z}{\mathbb{Z}}
% Symboles de validation
\providecommand{\cmark}{\ding{51}}
\providecommand{\xmark}{\ding{55}}
\providecommand{\wmark}{\ding{73}}
\newcommand{\symSucces}{\textcolor{success}{\cmark}}
\newcommand{\symEchec}{\textcolor{failure}{\xmark}}
\newcommand{\symAttention}{\textcolor{warning}{\wmark}}
% ==================== TABLE DES MATIÈRES ====================
\usepackage{tocloft}
\setlength{\cftbeforesecskip}{2pt}
\setlength{\cftsecindent}{0pt}
\setlength{\cftsubsecindent}{10pt}
\setlength{\cftsubsubsecindent}{20pt}
\setcounter{tocdepth}{3}   % profondeur de la table des matières
% ==================== GÉOMÉTRIE ET MISE EN PAGE ====================
\usepackage{geometry}
\geometry{
margin=2.5cm,
headsep=4mm,
footskip=12pt
}
% Interlignes et espacements
\setlength{\baselineskip}{14pt}
\setlength{\parskip}{6pt}
\setlength{\parindent}{0pt}
\setlength{\parsep}{0pt}
\setlength{\itemsep}{2pt}
\setlength{\topsep}{2pt}
\setlength{\abovedisplayskip}{8pt}
\setlength{\belowdisplayskip}{8pt}
\setlength{\abovedisplayshortskip}{4pt}
\setlength{\belowdisplayshortskip}{4pt}
% ==================== TABLEAUX AMÉLIORÉS ====================
\usepackage{array}
\usepackage{booktabs}
\usepackage{longtable}
\usepackage{colortbl}
\usepackage{multirow}
\usepackage{float}
\usepackage{adjustbox}
\setlength{\LTpre}{6pt}
\setlength{\LTpost}{6pt}
\setlength{\extrarowheight}{2pt}
\renewcommand{\arraystretch}{1.2}
% ==================== LISTES ====================
\usepackage{enumitem}
\setlist{nosep, left=0pt}
% ==================== COULEURS ====================
\usepackage{xcolor}
\definecolor{success}{RGB}{0,150,0}
\definecolor{failure}{RGB}{200,0,0}
\definecolor{warning}{RGB}{255,140,0}
\definecolor{codebg}{RGB}{245,245,245}
% ==================== CODE ET LISTINGS ====================
\usepackage{listings}
\usepackage{framed}
%\usepackage{fancyvrb}   % optionnel
\definecolor{shadecolor}{rgb}{0.9,0.9,0.9}
% Les environnements Shaded et Highlighting seront définis par pandoc via 
\lstset{
basicstyle=\ttfamily\small,
breaklines=true,
frame=single,
backgroundcolor=\color{codebg},
numbers=left,
numberstyle=\tiny\color{gray},
keywordstyle=\color{blue},
commentstyle=\color{green!50!black},
stringstyle=\color{red},
showstringspaces=false,
tabsize=4,
language=Python
}
% ==================== BOÎTES ET ENCADRÉS ====================
\usepackage{tikz}
\usepackage[framemethod=tikz]{mdframed}
\usepackage{tcolorbox}
\tcbuselibrary{most}
\newtcolorbox{validationbox}[2][]{
colback=green!5!white,
colframe=green!75!black,
fonttitle=\bfseries,
title=#2,
#1
}
\newtcolorbox{errorbox}[2][]{
colback=red!5!white,
colframe=red!75!black,
fonttitle=\bfseries,
title=#2,
#1
}
\newtcolorbox{warningbox}[2][]{
colback=orange!5!white,
colframe=orange!75!black,
fonttitle=\bfseries,
title=#2,
#1
}
% ==================== FIGURES ET IMAGES ====================
\usepackage{graphicx}
\usepackage{adjustbox}
\usepackage{capt-of}
\usepackage{pdflscape}

% Paramètres par défaut
\setkeys{Gin}{keepaspectratio=true}

% Commande figure auto-cadrée
\newcommand{\autofigure}[3][]{
\begin{figure}[htbp]
    \centering
    \adjustbox{max width=\textwidth, max height=0.6\textheight, keepaspectratio, center}{
        \includegraphics{#2}
    }
    \captionof{figure}{#3}
    \label{fig:#2}
\end{figure}
}

% Commande figure paysage
\newcommand{\autolandscapefigure}[3][]{
\afterpage{
\begin{landscape}
\begin{figure}[p]
    \centering
    \adjustbox{max width=0.9\textwidth, max height=0.7\textheight, keepaspectratio, center}{
        \includegraphics{#2}
    }
    \captionof{figure}{#3}
    \label{fig:#2}
\end{figure}
\end{landscape}
}
}
% ==================== EN-TÊTE ET PIED DE PAGE ====================
\usepackage{fancyhdr}
% Style normal
\pagestyle{fancy}
\fancyhf{}
\fancyhead[C]{\textsc{Modèle Cl(6,6) — Rapport de Synthèse}}
\fancyfoot[C]{\thepage}
\renewcommand{\headrulewidth}{0.4pt}
\renewcommand{\footrulewidth}{0pt}
% Style "plain" = même numérotation que fancy
\fancypagestyle{plain}{
\fancyhf{}
\fancyhead[C]{\textsc{Modèle Cl(6,6) — Rapport de Synthèse}}
\fancyfoot[C]{\thepage}
\renewcommand{\headrulewidth}{0.4pt}
\renewcommand{\footrulewidth}{0pt}
}
% Première page sans numéro
\fancypagestyle{firstpage}{
\fancyhf{}
\fancyhead[C]{\textsc{Modèle Cl(6,6)}}
\fancyfoot[C]{}
\renewcommand{\headrulewidth}{0.4pt}
\renewcommand{\footrulewidth}{0pt}
}
% ==================== HYPERLIENS ====================
\usepackage{hyperref}
\hypersetup{
colorlinks=true,
linkcolor=blue,
urlcolor=blue,
citecolor=blue
}
% ==================== PROTECTION LIGNES ORPHELINES ====================
\clubpenalty=10000
\widowpenalty=10000
\displaywidowpenalty=10000
% ==================== DÉFINITIONS POUR PANDOC ====================
\providecommand{\tightlist}{%
\setlength{\itemsep}{2pt}\setlength{\parskip}{6pt}
}
\newcommand{\passthrough}[1]{#1}
% ==================== VARIABLES PANDOC ====================
\title{``ANGULAR TRANSITIONS BETWEEN PARTICLES''}
\author{``Bruno DE DOMINICIS''}
\date{``March 2026''}
% ==================== REDÉFINITION DE \maketitle ====================
\makeatletter
\renewcommand{\maketitle}{
\begin{center}
\vspace*{1cm}
{\fontsize{18}{24}\selectfont \textbf{\@title} \par}
\vspace{1cm}
{\Large
\begin{tabular}[t]{c}
\@author
\end{tabular}\par}
\vspace{1cm}
{\Large \@date \par}
\end{center}
\vspace{1cm}
}
\makeatother
% ==================== CORPS DU DOCUMENT ====================
\begin{document}
\maketitle
\thispagestyle{firstpage}
% --- Mesure de la hauteur totale des abstracts ---
\newlength{\abstractheight}
\setbox0=\vbox{
\begin{center}
\begin{minipage}{0.85\textwidth}
\noindent \textbf{Résumé :} We present a reformulation of transitions
between elementary particles (beta decay, e+e- annihilation, \(pp\)
nuclear fusion, muon decay, neutrino oscillations, pair production) as
angle rearrangements in the Ibozoo Uû lattice, where each particle is
described by a pentad of the Clifford algebra \(\text{Cl}(6,6)\). The
generators \(i,j,k,I,J,K\) are interpreted as angles. We define a
transition operator \(T\) acting on the Hilbert space of pentads,
decomposed into structure, fire, water, and mixed parts, and we
establish selection rules (conservation of generators, chirality). An
angular formulation relates \(T\) to rotations in 6-dimensional space.
The example of \(\beta\) decay is discussed in detail. This formalism
offers a unified geometric view of interactions, in accordance with
Ummite concepts.
\end{minipage}
\end{center}
\vspace{0.5cm}
\begin{center}
\begin{minipage}{0.85\textwidth}
\noindent \textbf{Abstract:} We present a reformulation of transitions
between elementary particles (\(\beta\) decay, \(e^+e^-\) annihilation,
\(pp\) nuclear fusion, muon decay, neutrino oscillations, pair
production) as angular rearrangements within the Ibozoo Uû network,
where each particle is described by a pentad of the Clifford algebra
\(\text{Cl}(6,6)\). The generators \(i,j,k,I,J,K\) are interpreted as
angles. We define a transition operator \(T\) acting on the Hilbert
space of pentads, decomposed into structure, fire, water, and mixed
parts, and we establish selection rules (conservation of generators,
chirality). An angular formulation relates \(T\) to rotations in
6-dimensional space. The example of \(\beta\) decay is treated in
detail. This formalism provides a unified geometric view of
interactions, in agreement with Ummite concepts.
\end{minipage}
\end{center}
\vspace{1cm}
}
\abstractheight=\ht0
\advance\abstractheight by \dp0
% Calcul de l'espace restant sur la page après le titre
\newlength{\remaining}
\remaining=\dimexpr\textheight - \pagetotal\relax
% Si la hauteur des abstracts dépasse l'espace restant, on agrandit la page
\ifdim\abstractheight > \remaining
\enlargethispage{\dimexpr\abstractheight - \remaining\relax}
\fi
\begin{center}
\begin{minipage}{0.85\textwidth}
\noindent \textbf{Résumé :} We present a reformulation of transitions
between elementary particles (beta decay, e+e- annihilation, \(pp\)
nuclear fusion, muon decay, neutrino oscillations, pair production) as
angle rearrangements in the Ibozoo Uû lattice, where each particle is
described by a pentad of the Clifford algebra \(\text{Cl}(6,6)\). The
generators \(i,j,k,I,J,K\) are interpreted as angles. We define a
transition operator \(T\) acting on the Hilbert space of pentads,
decomposed into structure, fire, water, and mixed parts, and we
establish selection rules (conservation of generators, chirality). An
angular formulation relates \(T\) to rotations in 6-dimensional space.
The example of \(\beta\) decay is discussed in detail. This formalism
offers a unified geometric view of interactions, in accordance with
Ummite concepts.
\end{minipage}
\end{center}
\vspace{0.5cm}
\begin{center}
\begin{minipage}{0.85\textwidth}
\noindent \textbf{Abstract:} We present a reformulation of transitions
between elementary particles (\(\beta\) decay, \(e^+e^-\) annihilation,
\(pp\) nuclear fusion, muon decay, neutrino oscillations, pair
production) as angular rearrangements within the Ibozoo Uû network,
where each particle is described by a pentad of the Clifford algebra
\(\text{Cl}(6,6)\). The generators \(i,j,k,I,J,K\) are interpreted as
angles. We define a transition operator \(T\) acting on the Hilbert
space of pentads, decomposed into structure, fire, water, and mixed
parts, and we establish selection rules (conservation of generators,
chirality). An angular formulation relates \(T\) to rotations in
6-dimensional space. The example of \(\beta\) decay is treated in
detail. This formalism provides a unified geometric view of
interactions, in agreement with Ummite concepts.
\end{minipage}
\end{center}
\vspace{1cm}
% === AJOUT DE LA TABLE DES MATIÈRES ===
\tableofcontents
\newpage
% ======================================
\clearpage
\pagestyle{fancy}
\hypertarget{angular-transitions-between-particles}{%
\section{ANGULAR TRANSITIONS BETWEEN
PARTICLES}\label{angular-transitions-between-particles}}

\hypertarget{note-on-the-pre-geometric-framework}{%
\subsection{Note on the pre-geometric
framework}\label{note-on-the-pre-geometric-framework}}

Before addressing transitions between particles, it is essential to
define the framework in which they occur. Following approaches that seek
algebraic or geometric foundations for fundamental physics {[}1--4{]},
we adopt as a working hypothesis the idea that elementary particles can
be described as stable configurations within a pre-geometric algebraic
substrate.

In this approach, the fundamental degrees of freedom are not quantum
fields propagating over a fixed spatiotemporal background, but rather
angular relations between the generators of a Clifford algebra. A useful
analogy is to consider the substrate as a \textbf{``bundle of oriented
axes''} whose mutual orientations encode physical information.

It is essential to note that an isolated generator has no direct
physical meaning; it is the relational structure---\textbf{the
configuration of the angles between the generators}---that defines
observable quantities such as charge, mass, or flavor.

This perspective is motivated by several considerations:

\begin{itemize}
\tightlist
\item
  Clifford algebras naturally encode spin, chirality, and gauge
  structures {[}1, 2{]};
\item
  Pre-geometric approaches suggest that spacetime itself might emerge
  from more fundamental combinatorial or algebraic relations {[}3, 4{]};
\item
  Representing particles as algebraic ``patterns'' offers a path to
  unify interactions through geometric rearrangements rather than
  through force exchanges.
\end{itemize}

\textbf{Link to our pentad model}: In this work, we use this concept to
construct a geometric representation of particles. A particle is not an
entity embedded in this lattice; it is a local and stable configuration
of this lattice.

Our formalism in Cl(6,6) and our \textbf{pentads} constitute a
mathematical model of this idea: 1. The six \textbf{generators}
\(i,j,k,I,J,K\) represent the fundamental ``axes'' of this algebraic
substrate in our cosmos. 2. The \textbf{angles} mentioned throughout the
text (for example, \(\theta_1\theta_4\)) are the mathematical
translation of the angular relationships between these axes within the
substrate. 3. A \textbf{pentad}, this set of five elements (three
structural bivectors, a ``fire'' axial element, a ``water'' polar
element), is the complete angular signature that defines the quantum
state and nature of a particle (its mass, charge, flavor) as a specific
and coherent pattern of these angular relationships within the lattice.

Thus, when we describe a transition \(A \to B + C + \dots\) as a
\textbf{``rearrangement of angles} within the \textbf{Fundamental
Angular Lattice}, we mean that the configuration of angular relations
that defined particle \(A\) dissolves and reconfigures to form new
stable configurations, those of particles \(B\) and \(C\). Particle
physics thus becomes a \textbf{dynamic geometry of angles}.

\begin{center}\rule{0.5\linewidth}{0.5pt}\end{center}

\hypertarget{principle-of-pentad-coding}{%
\subsection{Principle of pentad
coding}\label{principle-of-pentad-coding}}

In our model, each elementary particle is represented by a
\textbf{pentad}, that is, an ordered set of five elements of the
Clifford algebra \(\text{Cl}(6,6)\). These elements are interpreted as
\textbf{angles} in a six-dimensional space (generators \(i,j,k,I,J,K\)).

They are divided into three roles:

\begin{itemize}
\tightlist
\item
  \textbf{The structure}: three bivectors (products of two generators)
  that define the particle's intrinsic identity. For nucleons, we choose
  \$\{iI, jJ, kK\}
  \(. This choice reflects the symmetry between the three spatial dimensions (\)i,j,k\() and the three charge dimensions (\)I,J,K\$).
  It is common to both the neutron and the proton because they share the
  same valence quark composition (two up quarks and one down quark for
  the proton; one up and two down for the neutron) but with different
  charges.
\item
  \textbf{Fire}: an element of the form \(i'v\), where \(i'\) is an
  additional generator (often associated with chirality or the weak
  interaction) and \(v\) is one of the six generators. For nucleons, we
  have \(i'k\), which represents the weak component. This element is
  identical for the neutron and the proton because the weak interaction
  acts similarly on both (they can transform into one another).
\item
  \textbf{Water}: an element of the form \(1v\), where \(1\) is a scalar
  element (or a neutral generator) and \(v\) is a generator,
  representing the \textbf{mass} or the \textbf{substance} along a
  preferred axis. It is this element that accounts for the difference
  between the neutron and the proton. The electric charge is encoded by
  the orientation of this axis: \(1i\) corresponds to a mass aligned
  along \(i\) (neutron, zero charge), while \(1j\) corresponds to an
  alignment along \(j\) (proton, positive charge). This interpretation
  is consistent with the idea that charge is a directional property in
  angle space.
\end{itemize}

\hypertarget{ux3b2-decay-n-p-e-ux3bdux2091}{%
\subsection{β Decay: n → p + e⁻ +
ν̄ₑ}\label{ux3b2-decay-n-p-e-ux3bdux2091}}

Thus, β⁻ decay is viewed as an \textbf{angular rearrangement}: the
neutron changes its substance angle from \(i\) to \(j\), which is
accompanied by the emission of the electron (which has \(1j\)) and the
antineutrino (which has \(1k\)). The other angles (structure and fire)
recombine to form the pentads of the final particles.

\hypertarget{initial-and-final-pentads}{%
\subsubsection{Initial and Final
Pentads}\label{initial-and-final-pentads}}

\textbf{Initial State} (neutron):

\[
P(n) = \{ iI,\ jJ,\ kK,\ i'k,\ 1i \}
\]

Associated angles (with the correspondence
\(i \leftrightarrow \theta_1\), \(j \leftrightarrow \theta_2\),
\(k \leftrightarrow \theta_3\), \(I \leftrightarrow \theta_4\),
\(J \leftrightarrow \theta_5\), \(K \leftrightarrow \theta_6\)): \[
\{\theta_1\theta_4,\ \theta_2\theta_5,\ \theta_3\theta_6,\ (\theta_1\theta_2\theta_3\theta_4\theta_5\theta_6)\cdot\theta_3,\ 1\cdot\theta_1\}
\]

\textbf{Final state} (proton + electron + antineutrino):

\[
P(p) = \{ iI,\ jJ,\ kK,\ i'k,\ 1j \}
\] \[
P(e^-) = \{ iI,\ iJ,\ iK,\ i'k,\ 1j \}
\] \[
P(\bar{\nu}_e) = \{ iK,\ iJ,\ iI,\ i'j,\ 1k \}
\]

(Note: The representation of the antineutrino is taken here to be
identical to that of the electron neutrino used in the section on
oscillations; a more precise convention could introduce signs to
distinguish between particle and antiparticle, but this does not affect
the qualitative discussion.)

\hypertarget{description-of-the-angular-transition}{%
\subsubsection{Description of the angular
transition}\label{description-of-the-angular-transition}}

\begin{itemize}
\tightlist
\item
  \textbf{Step 1}: The neutron changes its substance angle from \(1i\)
  to \(1j\) (rotation of the mass axis from \(i\) to \(j\)). This is the
  transformation of the down quark into an up quark via the weak
  interaction.
\item
  \textbf{Step 2}: The neutron's structure angles \(\{iI, jJ, kK\}\)
  reorganize to yield those of the proton (unchanged) and those of the
  electron (which are \(\{iI, iJ, iK\}\)). The electron carries away
  part of the structure.
\item
  \textbf{Step 3}: The antineutrino is formed from the remaining angles,
  with a \(1k\) component and an \(i'j\) component (i.e., a permutation
  of the generators).
\end{itemize}

This geometric description advantageously replaces the usual picture of
the emission of a \(W^-\) boson: the transition is a pure rearrangement
of angles, without the introduction of virtual particles.

\begin{center}\rule{0.5\linewidth}{0.5pt}\end{center}

\hypertarget{annihilation-e-e-ux3b3ux3b3}{%
\subsection{Annihilation e⁺ e⁻ → γγ}\label{annihilation-e-e-ux3b3ux3b3}}

\hypertarget{principle-of-pentad-coding-for-leptons-and-photons}{%
\subsubsection{Principle of pentad coding for leptons and
photons}\label{principle-of-pentad-coding-for-leptons-and-photons}}

\begin{itemize}
\tightlist
\item
  \textbf{The structure}: for charged leptons, \(\{iI, iJ, iK\}\); for
  the photon, the same but without fire or water.
\item
  \textbf{Fire}: \(i'k\) for the electron, \(-i'k\) for the positron.
\item
  \textbf{Water}: \(1j\) for the electron, \(-1j\) for the positron.
\end{itemize}

Photons are represented by \(\{iI, iJ, iK, 0, 0\}\).

\hypertarget{initial-and-final-pentads-1}{%
\subsubsection{Initial and final
pentads}\label{initial-and-final-pentads-1}}

\textbf{Initial state}:

\[
P(e^-) = \{ iI,\ iJ,\ iK,\ i'k,\ 1j \}
\] \[
P(e^+) = \{ -iI,\ -iJ,\ -iK,\ -i'k,\ -1j \}
\]

\textbf{Final state} (two photons):

\[
P(\gamma_1) = \{ iI,\ iJ,\ iK,\ 0,\ 0 \}
\] \[
P(\gamma_2) = \{ iI,\ iJ,\ iK,\ 0,\ 0 \}
\]

\hypertarget{description-of-the-angular-transition-1}{%
\subsubsection{Description of the angular
transition}\label{description-of-the-angular-transition-1}}

Annihilation results from a \textbf{cancellation of opposite angles}:
each component of the positron is the opposite of that of the electron.
Their sum is zero, but energy must be conserved. The six bivectors
(three of each) recombine into two sets of three, forming two photons.
The fire and water components actually cancel each other out, since
photons have neither mass nor weak interaction.

In terms of angles:

\[
\{iI,iJ,iK,i'k,1j\} + \{-iI,-iJ,-iK,-i'k,-1j\} \to \{iI,iJ,iK,0,0\} + \{iI,iJ,iK,0,0\}
\]

The conservation of the number of generators is respected: each photon
has three bivectors (six generators), while the pair also had six (with
opposite signs).

\begin{center}\rule{0.5\linewidth}{0.5pt}\end{center}

\hypertarget{nuclear-fusion-p-p-d-e-ux3bdux2091}{%
\subsection{Nuclear fusion: p + p → d + e⁺ +
νₑ}\label{nuclear-fusion-p-p-d-e-ux3bdux2091}}

\hypertarget{pentads}{%
\subsubsection{Pentads}\label{pentads}}

\textbf{Initial state} (two protons):

\[
P(p_1) = \{ iI,\ jJ,\ kK,\ i'k,\ 1j \}
\] \[
P(p_2) = \{ iI,\ jJ,\ kK,\ i'k,\ 1j \}
\]

\textbf{Final state} (deuteron, positron, electron neutrino):

\[
P(d) = \{ iI,\ jJ,\ kK,\ i'k,\ 1j \} \quad (\text{with binding})
\] \[
P(e^+) = \{ -iI,\ -iJ,\ -iK,\ -i'k,\ -1j \}
\] \[
P(\nu_e) = \{ iK,\ iJ,\ iI,\ i'j,\ 1k \}
\]

\hypertarget{description}{%
\subsubsection{Description}\label{description}}

One of the protons undergoes an internal β⁺ decay: its spin changes from
\(1j\) to \(1i\) (thus becoming a neutron), and it emits a positron
(opposite configuration) and a neutrino (permutated structure) . The
neutron thus formed binds with the other proton to form a deuteron. The
binding is represented by an additional coupling of the angles, not
detailed in the pentad itself.

The conservation of the generators is ensured by the selection rules of
the weak interaction, which allow for flavor changes.

\begin{center}\rule{0.5\linewidth}{0.5pt}\end{center}

\hypertarget{muon-decay-ux3bc-e-ux3bdux2091-ux3bdux3bc}{%
\subsection{Muon decay: μ⁻ → e⁻ + ν̄ₑ +
νμ}\label{muon-decay-ux3bc-e-ux3bdux2091-ux3bdux3bc}}

\hypertarget{pentads-1}{%
\subsubsection{Pentads}\label{pentads-1}}

\textbf{Muon}:

\[
P(\mu^-) = \{ jI,\ jJ,\ jK,\ i'i,\ 1k \}
\]

\textbf{Electron}:

\[
P(e^-) = \{ iI,\ iJ,\ iK,\ i'k,\ 1j \}
\]

\textbf{Electron antineutrino}:

\[
P(\bar{\nu}_e) = \{ iK,\ iJ,\ iI,\ i'j,\ 1k \}
\]

\textbf{Muon neutrino}:

\[
P(\nu_\mu) = \{ jK,\ jI,\ jJ,\ i'k,\ 1i \}
\]

(Note: The representation of antineutrinos here is modeled after that of
neutrinos, which is a simplification. A more detailed description could
introduce signs to distinguish between particles and antiparticles, but
this does not change the geometric interpretation of the axis rotation.)

\hypertarget{description-1}{%
\subsubsection{Description}\label{description-1}}

The muon has its structure along the \(j\) axis, the electron along the
\(i\) axis. The decay corresponds to a \textbf{rotation of the flavor
axis} from \(j\) to \(i\), accompanied by the emission of two neutrinos
that carry the permutations. The fire and water elements redistribute
themselves among the three final products.

\begin{center}\rule{0.5\linewidth}{0.5pt}\end{center}

\hypertarget{neutrino-oscillation-ux3bdux2091-ux3bdux3bc}{%
\subsection{Neutrino oscillation: νₑ →
νμ}\label{neutrino-oscillation-ux3bdux2091-ux3bdux3bc}}

\hypertarget{pentads-2}{%
\subsubsection{Pentads}\label{pentads-2}}

\textbf{Electron neutrino}:

\[
P(\nu_e) = \{ iK,\ iJ,\ iI,\ i'j,\ 1k \}
\]

\textbf{Muon neutrino}:

\[
P(\nu_\mu) = \{ jK,\ jI,\ jJ,\ i'k,\ 1i \}
\]

\hypertarget{description-2}{%
\subsubsection{Description}\label{description-2}}

Oscillation is a \textbf{continuous rotation of the angles in space}:
the preferred axis of the structure gradually shifts from \(i\) to
\(j\). The transition operator is \(\exp(i\alpha L_{ji})\), where
\(\alpha\) is proportional to time and \(\Delta m^2\). The oscillation
probability is \(\sin^2\alpha\), consistent with the standard formula.

\begin{center}\rule{0.5\linewidth}{0.5pt}\end{center}

\hypertarget{pair-production-ux3b3-e-e}{%
\subsection{Pair production: γ → e⁺ +
e⁻}\label{pair-production-ux3b3-e-e}}

\hypertarget{pentads-3}{%
\subsubsection{Pentads}\label{pentads-3}}

\textbf{Photon} (in the presence of an external field):

\[
P(\gamma) = \{ iI,\ iJ,\ iK,\ 0,\ 0 \}
\]

\textbf{Electron-positron pair}:

\[
P(e^-) = \{ iI,\ iJ,\ iK,\ i'k,\ 1j \}
\] \[
P(e^+) = \{ -iI,\ -iJ,\ -iK,\ -i'k,\ -1j \}
\]

\hypertarget{description-3}{%
\subsubsection{Description}\label{description-3}}

A high-energy photon interacts with an external field (for example, the
Coulomb field of a nucleus). This field provides the missing
orientations (axes \(j\), \(k\), \(i'\)) that allow for the
\textbf{creation of matter}: the water elements \(1j\) and \(-1j\) as
well as the fire elements \(i'k\) and \(-i'k\) appear, while the
photon's structure splits into two opposing sets. This is the reverse of
annihilation.

\begin{center}\rule{0.5\linewidth}{0.5pt}\end{center}

\hypertarget{summary-table-of-transitions}{%
\subsection{Summary table of
transitions}\label{summary-table-of-transitions}}

\begin{longtable}[]{@{}
  >{\raggedright\arraybackslash}p{(\columnwidth - 4\tabcolsep) * \real{0.2857}}
  >{\raggedright\arraybackslash}p{(\columnwidth - 4\tabcolsep) * \real{0.1786}}
  >{\raggedright\arraybackslash}p{(\columnwidth - 4\tabcolsep) * \real{0.5357}}@{}}
\toprule
\begin{minipage}[b]{\linewidth}\raggedright
Transition
\end{minipage} & \begin{minipage}[b]{\linewidth}\raggedright
Type
\end{minipage} & \begin{minipage}[b]{\linewidth}\raggedright
Angular transformation
\end{minipage} \\
\midrule
\endhead
\(n \to p + e^- + \bar{\nu}_e\) & \(\beta^-\) & Rotation of the
substance angle (\(1i \to 1j\)) + emission \\
\(e^+e^- \to \gamma\gamma\) & Annihilation & Cancellation of opposite
angles \\
\(pp \to e^+ \nu_e\) & Fusion & Entanglement of angles (with proton
transformation) \\
\(\mu^- \to e^- \bar{\nu}_e\nu_\mu\) & Decay & Rotation of the flavor
axis (\(j \to i\)) \\
\(\nu_e \rightleftarrows \nu_\mu\) & Oscillation & Rotation of angles in
space (time-dependent) \\
\(\gamma \to e^+e^-\) & Pair production & Creation of matter (water and
fire) from angles, assisted by an external field \\
\bottomrule
\end{longtable}

\begin{center}\rule{0.5\linewidth}{0.5pt}\end{center}

\hypertarget{angular-conservation-laws}{%
\subsection{Angular Conservation Laws}\label{angular-conservation-laws}}

\begin{itemize}
\tightlist
\item
  \textbf{Conservation of the total number of generators}: each element
  (bivector, fire, water) has a fixed number of generators; the sum is
  invariant.
\item
  \textbf{Conservation of generator type}: \(i,j,k,I,J,K\) (and \(i'\))
  are conserved individually, modulo gauge transformations.
\item
  \textbf{Conservation of chirality}: linked to the presence of \(i'\)
  in the fire elements.
\item
  \textbf{Conservation of total angular momentum}: in angle space,
  rotations preserve the norm of the angular momentum vector.
\end{itemize}

\begin{center}\rule{0.5\linewidth}{0.5pt}\end{center}

\hypertarget{mathematical-formalization}{%
\subsection{Mathematical
formalization}\label{mathematical-formalization}}

Let \(\mathcal{H}\) be the Hilbert space of pentads (dimension 72 or
144). A transition operator \(T\) acts on tensor products:

\[
\langle P_B \otimes P_C \otimes \cdots | T | P_A \rangle
\]

The transition probability is given by the square of the matrix element.

\begin{center}\rule{0.5\linewidth}{0.5pt}\end{center}

\hypertarget{pentadic-transition-operator}{%
\section{PENTADIC TRANSITION
OPERATOR}\label{pentadic-transition-operator}}

\hypertarget{decomposition}{%
\subsection{Decomposition}\label{decomposition}}

\[
T = T_{\text{structure}} + T_{\text{fire}} + T_{\text{water}} + T_{\text{mixed}}
\]

\begin{itemize}
\tightlist
\item
  \(T_{\text{structure}}\) acts on triples of bivectors.
\item
  \(T_{\text{fire}}\) acts on the elements \(i'v\).
\item
  \(T_{\text{water}}\) acts on the elements \(1v\).
\item
  \(T_{\text{mixed}}\) pairs the different types.
\end{itemize}

\hypertarget{compact-expression}{%
\subsection{Compact Expression}\label{compact-expression}}

In \(\text{Cl}(6,0)\), with a basis \(\{E_a\}\):

\[
T = \sum_{a,b} g_{ab} (E_a \otimes E_b^*)
\]

For a pentad \(P = \{B_1,B_2,B_3,F,S\}\):

\[
|P\rangle = |B_1\rangle \wedge |B_2\rangle \wedge |B_3\rangle \wedge |F\rangle \wedge |S\rangle
\]

\[
T|P\rangle = \sum_{P'} T_{P'P} |P'\rangle,
\]

\[
T_{P'P} = \sum_{\text{cycles}} \prod_{k=1}^5 \langle e'_k | T_{\text{local}} | e_k \rangle
\]

\hypertarget{selection-rules}{%
\subsection{Selection rules}\label{selection-rules}}

\begin{itemize}
\tightlist
\item
  Conservation of the number and type of generators.
\item
  Conservation of chirality.
\item
  Pentadic Feynman diagrams.
\end{itemize}

\hypertarget{example-ux3b2-decay}{%
\subsection{Example: β decay}\label{example-ux3b2-decay}}

\[
\mathcal{M} = \langle P(p) \otimes P(e^-) \otimes P(\bar{\nu}_e) | T | P(n) \rangle = G_F \, J_{\text{had}} \cdot J_{\text{lep}}
\]

where

\[
J_{\text{had}} = \langle P(p) | T_{\text{struct}} \otimes T_{\text{fire}} | P(n) \rangle,
\]

\[
J_{\text{lep}} = \langle P(e^-) \otimes P(\bar{\nu}_e) | T_{\text{fire}} \otimes T_{\text{water}} |0\rangle.
\]

\hypertarget{canonical-quantization}{%
\subsection{Canonical quantization}\label{canonical-quantization}}

Operators \(a_P\), \(a_P^\dagger\) with
\([a_P, a_Q^\dagger]_{\mp} = \delta_{PQ}\). The interaction Hamiltonian
is a sum of products.

\hypertarget{angular-formulation}{%
\subsection{Angular formulation}\label{angular-formulation}}

The angles \(\theta_\mu\) are the generators. \(T\) becomes a
differential operator:

\[
T = \exp\!\left( i \sum_{i,j} \omega_{ij} L_{ij} \right), \quad L_{ij} = -i(\theta_i \partial_{\theta_j} - \theta_j \partial_{\theta_i})
\]

For the decay \(\mu \to e\), the probability is \(\sin^2\alpha\), with
\(\alpha \propto t \Delta m^2/E\).

\begin{center}\rule{0.5\linewidth}{0.5pt}\end{center}

\hypertarget{toward-a-complete-theory}{%
\subsection{Toward a complete theory}\label{toward-a-complete-theory}}

It remains to: - Determine the coupling constants from the data. -
Derive the pentadic Feynman rules. - Calculate the cross sections. -
Predict new processes.

\hypertarget{conclusion}{%
\subsection{Conclusion}\label{conclusion}}

The operator \(T\), defined by a sum over cycles of local matrix
elements, provides a coherent mathematical framework for describing all
particle transitions as angle rearrangements. This geometric model,
inspired by Ummite concepts, unifies interactions and paves the way for
a deeper understanding of particle physics.
\end{document}
